% Part: set-theory
% Chapter: story
% Section: extensionality

\documentclass[../../../include/open-logic-section]{subfiles}

\begin{document}

\olfileid{sth}{story}{extensionality}
\olsection{সদস্যভিত্তিক সমতা}

প্রথমেই বলতে হয়, একটি সেটের পরিচয় তার !!{element}s দিয়েই
নির্ধারিত হয়। আরও নির্দিষ্টভাবে:

\begin{axiom}[সদস্যভিত্তিক সমতা]
যদি সেট $A$ ও~$B$-এর !!{element}s একই হয়, তবে $A$ ও~$B$
একই সেট।
\[
  \lforall[A][\lforall[B][(\lforall[x][(x \in A \liff x \in B)] \lif
  \eq[A][B])]]
\]
\end{axiom}

\olref[sfr][][]{part}-এর সর্বত্র আমরা এই নীতি ধরে নিয়েছি। তবু
কথাটি আবার বলা দরকার। সদস্যভিত্তিক সমতার স্বতঃসিদ্ধটি এই মৌলিক
ধারণাই প্রকাশ করে যে, একটি সেট তার !!{element}s দ্বারা নির্ধারিত।
(এই দিক থেকে সেটের সঙ্গে \emph{ধারণা}-র তুলনা করা যায়: ঠিক একই
বস্তু বহু ভিন্ন ধারণার অধীনে পড়তে পারে।)

এই নীতি গ্রহণ করব কেন? বলা যুক্তিসঙ্গত যে, সদস্যভিত্তিক সমতা
অস্বীকার করলে এমন কিছুকেই পরিত্যাগ করা হয়, যাকে আদৌ
\emph{সেটতত্ত্ব} বলা চলে। সেটতত্ত্ব আসলে সদস্যভিত্তিক সংগ্রহেরই
তত্ত্ব; এর কমও নয়, বেশিও নয়।

কিন্তু \olref[sth][][]{part}-এ প্রকৃত চ্যালেঞ্জ হলো এমন নীতি
নির্ধারণ করা, যা বলে দেয় \emph{কোন কোন সেট আছে}। দেখা যাবে,
এই প্রশ্নের একমাত্র আপাত ``সুস্পষ্ট'' উত্তরটি প্রমাণসাপেক্ষে ভুল।

\end{document}
